\documentclass[UTF8]{ctexart}
\usepackage[a4paper,margin=2.4cm]{geometry}
\usepackage{hyperref}
\usepackage{longtable}
\usepackage{enumitem}
\usepackage{xcolor}

\hypersetup{
  colorlinks=true,
  linkcolor=blue,
  urlcolor=blue
}

\setlist[itemize]{leftmargin=1.8em}
\setlist[enumerate]{leftmargin=1.8em}

\title{TE-KG 智能体系统架构与修复说明}
\author{TE-KG Project}
\date{\today}

\begin{document}
\maketitle

\section{文档目的}

本文档记录 TE-KG 当前智能体系统的运行方式，以及本轮针对明显稳定性问题实施的修复。
范围仅包括 Deep Think、Agent、插件执行、LLM 调用、异步运行编排和前端智能体体验。
不覆盖首页可视化、BioRender、TE taxonomy UI 或普通内容页样式。

\section{当前模式划分}

\begin{itemize}
  \item Deep Think 是轻量模式：单模型工具增强推理，不显示六阶段工作流，适合序列、关系、分类、表达和基因组位置等直接查询。
  \item Agent 是重型模式：异步 run 加 worker，显示 Understanding、Planning、Collecting、Executing、Integrating、Writing 六个阶段，适合机制、文献证据、图谱分析和多步综合。
  \item 前端入口为 \texttt{agent.php} 和 \texttt{assets/js/pages/agent.js}。Deep Think 调用 \texttt{api/deep\_think\_stream.php}；Agent 创建 run 调用 \texttt{api/agent\_runs.php}，轮询状态调用 \texttt{api/agent\_run\_status.php}。
\end{itemize}

\section{本轮修复}

\begin{longtable}{p{0.22\linewidth}p{0.70\linewidth}}
\textbf{问题} & \textbf{修复方式} \\
\hline
Writing 超时配置未生效 & 移除 Agent 和 Deep Think 中对 \texttt{llm\_answer\_chat\_timeout} 的 20 秒硬上限，使 \texttt{config.local.php} 中的 40 秒配置真实生效；reasoner timeout 也尊重配置值。 \\
位置问题错路由 & 在实体归一化和意图识别中补充 ``位于哪里''、``在哪里''、``定位''、``基因组位置''、``染色体位置'' 等 genome/location 触发词；expression 判断仍早于 genome，避免表达问题被误判。 \\
sufficiency 超时后过早停止 & 当 sufficiency LLM 无可用结果且非 hard-stop intent 仍有 primary path 未执行插件时，继续推荐剩余 primary path，不再直接判定证据充分。 \\
Agent 过程旁白消耗 LLM & Agent 常规 analysis、planning、tool selected、tool result、reflection、synthesizing 事件改为模板优先，减少 narrator LLM 调用和 6 秒级超时噪声。 \\
陈旧 run 不会失败 & 在状态查询层增加 stale watchdog：pending 超过 30 秒未启动标记失败；running 超过 \texttt{agent\_execution\_timeout + 30} 秒标记失败。 \\
\end{longtable}

\section{插件职责}

\begin{itemize}
  \item Entity Resolver：稳定实体名称、别名链和 broad alias 边界。
  \item Graph Plugin：查询本地图谱中的实体邻域和结构化关系。
  \item Graph Analytics Plugin：执行全局排名、计数、关联度和拓扑类查询。
  \item Cypher Explorer Plugin：在固定模板不足时生成只读 Cypher 探索查询。
  \item Sequence、Genome、Expression、Tree Plugin：分别提供序列、基因组位置、表达和分类谱系信息。
  \item Literature、Literature Reading、Citation Resolver：检索文献、综合论文证据、归一化引用。
\end{itemize}

\section{Mermaid 流程图源码}

下面的 Mermaid 代码按人工语法检查规则编写：节点 ID 使用 ASCII，标签使用双引号，判断节点使用花括号，所有 subgraph 均闭合，class 定义位于末尾。

\begin{verbatim}
flowchart TD
  U["User question on agent.php"] --> JS["agent.js reads TEKG_AGENT_CONFIG"]
  JS --> M{"Selected mode"}
  M -->|Deep Think| DT_API["POST api/deep_think_stream.php"]
  M -->|Agent| AR_CREATE["POST api/agent_runs.php"]

  subgraph FE["Frontend runtime"]
    JS --> DT_UI["Deep thinking panel without six stage graph"]
    JS --> AG_UI["Agent six stage workflow"]
    AG_UI --> POLL["Poll api/agent_run_status.php with after sequence"]
    POLL --> RENDER["Render stage state, tool cards, answer, or failure"]
  end

  subgraph CFG["Runtime configuration"]
    LOCAL["api/config.local.php"] --> BOOT["tekg_agent_config"]
    BOOT --> MODELS["deepseek-v4-flash models and timeout values"]
    MODELS --> LLM["TekgAgentLlmClient"]
    LLM --> RELAY{"llm_relay_url configured"}
    RELAY -->|yes| RELAY_CALL["POST relay chat endpoint"]
    RELAY -->|no| PROVIDER_CALL["Direct provider call"]
  end

  subgraph DT["Deep Think path"]
    DT_API --> DT_SVC["DeepThinkService"]
    DT_SVC --> DT_ANALYZE["EntityNormalizer analyzes intent and entities"]
    DT_ANALYZE --> DT_ENTITY["Entity Resolver"]
    DT_ENTITY --> DT_ROUTE{"Simple deterministic route"}
    DT_ROUTE -->|sequence| DT_SEQ["Sequence Plugin"]
    DT_ROUTE -->|relationship| DT_GRAPH["Graph Plugin"]
    DT_ROUTE -->|classification| DT_TREE["Tree Plugin"]
    DT_ROUTE -->|expression| DT_EXPR["Expression Plugin"]
    DT_ROUTE -->|genome| DT_GENOME["Genome Plugin"]
    DT_ROUTE -->|mechanism or papers| DT_LIT["Graph plus Literature plus Reading"]
    DT_SEQ --> DT_SYN["Aggregate evidence and citations"]
    DT_GRAPH --> DT_SYN
    DT_TREE --> DT_SYN
    DT_EXPR --> DT_SYN
    DT_GENOME --> DT_SYN
    DT_LIT --> DT_SYN
    DT_SYN --> DT_DET{"Deterministic answer available"}
    DT_DET -->|yes| DT_DIRECT["Render full sequence or full relation list"]
    DT_DET -->|no| DT_WRITE["LLM direct answer"]
    DT_DIRECT --> DT_DONE["SSE answer event"]
    DT_WRITE --> DT_DONE
  end

  subgraph AG["Async Agent path"]
    AR_CREATE --> SAVE["Save payload.json and state.json"]
    SAVE --> KICK["Local kickoff request or CLI worker"]
    KICK --> WORKER["agent_run_execute marks run running"]
    WORKER --> AG_SVC["AcademicAgentService"]
    AG_SVC --> UNDER["Understanding stage"]
    UNDER --> PLAN["Planning stage"]
    PLAN --> COLLECT["Collecting stage builds plugin queue"]
    COLLECT --> EXEC["Executing stage runs next plugin"]
    EXEC --> PLUGINS{"Plugin selected"}
    PLUGINS --> P_ENTITY["Entity Resolver"]
    PLUGINS --> P_GRAPH["Graph Plugin"]
    PLUGINS --> P_ANALYTICS["Graph Analytics Plugin"]
    PLUGINS --> P_CYPHER["Cypher Explorer Plugin"]
    PLUGINS --> P_SEQ["Sequence Plugin"]
    PLUGINS --> P_GENOME["Genome Plugin"]
    PLUGINS --> P_EXPR["Expression Plugin"]
    PLUGINS --> P_TREE["Tree Plugin"]
    PLUGINS --> P_LIT["Literature Plugin"]
    PLUGINS --> P_READ["Literature Reading Plugin"]
    P_ENTITY --> SUFF["Minimum gate and sufficiency check"]
    P_GRAPH --> SUFF
    P_ANALYTICS --> SUFF
    P_CYPHER --> SUFF
    P_SEQ --> SUFF
    P_GENOME --> SUFF
    P_EXPR --> SUFF
    P_TREE --> SUFF
    P_LIT --> SUFF
    P_READ --> SUFF
    SUFF -->|need more evidence| COLLECT
    SUFF -->|enough evidence| CITE["Citation Resolver if citations exist"]
    CITE --> INTEGRATE["Integrating stage aggregates claims citations limits"]
    INTEGRATE --> STRUCT["answer_structure JSON or fallback"]
    STRUCT --> WRITE["Writing stage uses trimmed payload"]
    WRITE --> DONE_EVENT["done event updates run state"]
    DONE_EVENT --> POLL
  end

  subgraph FIX["Implemented fixes"]
    F1["Fix timeout helpers to honor config 40 seconds"]
    F2["Route location wording to Genome Plugin"]
    F3["On sufficiency LLM timeout continue remaining primary path"]
    F4["Use deterministic templates for routine Agent narration"]
    F5["Mark stale async runs failed in status endpoint"]
  end

  F1 --> WRITE
  F2 --> DT_ANALYZE
  F2 --> UNDER
  F3 --> SUFF
  F4 --> UNDER
  F4 --> PLAN
  F4 --> EXEC
  F5 --> POLL

  classDef fix fill:#fff7d6,stroke:#b7791f,color:#111;
  class F1,F2,F3,F4,F5 fix;
\end{verbatim}

\end{document}
